\subsection{The unit ordinary eighth moment for terminal $(1,\allowbreak 2,\allowbreak 3)$}\label{proof:Q5-C062}
\paragraph{Q62.1 Model, directions and local bound}

Let $(X,\allowbreak  \mu)$ be a nonempty finite probability space and let $K$ be real and skew with $|K(x,\allowbreak  y)|\le\allowbreak 1$.
All integrals, adjoints, norms and spectral measures use $\mu$. Put

\[
Jv=(Ev)1,\quad A=J+K,\quad T=A^*=J-K,
\quad g=K1,\ p=1+g,\ q=1-g,\ d=Eg^2.
\]

Write $f_n=\allowbreak A^n\one$, $\bar f_n=\allowbreak T^n\one$, $h_n=\allowbreak Tf_{n-1}$, and

\[
r=Kg,\quad s=K^2g,\quad \delta_j=\|K^j1\|_2^2,
\quad \theta=13/32.
\]

The finite calculus of Section \ref{q3:app:foundations} gives
$0\le\allowbreak  d\le\allowbreak 1/3$, $d^2\le\allowbreak  \delta_2\le\allowbreak \theta d$, $\|K\|^2\le\allowbreak \theta$,
$|Kv(x)|\le\allowbreak \|v\|_2$, path positivity and the equal-arm bound
$0\le\allowbreak  f_\ell h_\ell\le\allowbreak 1$.


The terminal residual lengths are $(1,\allowbreak  2,\allowbreak  3)$, with messages $q,\allowbreak  f_2,\allowbreak  h_3$, respectively.
The unit arm points toward the terminal, the two-edge arm points outward, and only the first edge of the three-edge arm points toward the terminal.
Allow any finite number of equal-arm pairs and put

\[
\begin{gathered}
E_0=\prod_j f_{\ell_j}h_{\ell_j},\qquad\\
F=E_0q f_2h_3,\qquad X=TF,\qquad H=pX.
\end{gathered}\tag{Q62.1.1}\label{eq:Q62:Q62:1:1}
\]

The connector has one edge, directed from the terminal to the receiving root, whose unit arm points outward.
These choices depend only on the arm lengths and pair multiplicities.

For every such structure and every finite host,

\[
\boxed{EH^8\le e^{-d/200}\le1-d/601.}
\tag{Q62.1.2}\label{eq:Q62:Q62:1:2}
\]

The bound also permits equal-arm decorations at the receiving root. Full reversal gives the corresponding
ordinary eighth-moment bound for $qAF^-$.



\paragraph{Q62.2 Positive envelope with the same kernel parameters}

By $E_0\le\allowbreak 1$ and positivity of $T$, it suffices to consider $P=\allowbreak qf_2h_3$.
Define

\[
t=1-d/2,\qquad k=t^{-1},\qquad a=t^2,
\qquad 5/6\le t\le1.
\]

The path identities are

\[
f_2=p+r,\quad h_3=q-d-r-s,\quad \bar f_3=q-d+r-s,
\quad f_2+h_3=2-d-s.
\]

Arithmetic--geometric mean for the two nonnegative factors gives

\[
0\le P\le q(t-s/2)^2.
\tag{Q62.2.1}\label{eq:Q62:Q62:2:1}
\]

To establish positivity of $H_0=\allowbreak q-ks$, use $h_3,\allowbreak  \bar f_3\ge\allowbreak 0$, which imply
$q-s\ge\allowbreak  d+|r|$, together with
$|s|\le\allowbreak \sqrt{\delta_2}\le\allowbreak \sqrt{13/96}<\allowbreak 3/8$. Thus

\[
H_0=(q-s)-(k-1)s
\ge d-\frac{d}{2t}|s|\ge31d/40\ge0.
\]

Since $Es=\allowbreak 0$, we have $EH_0=\allowbreak 1$. The kernel of $T$ lies in $[0,\allowbreak  2]$, so

\[
B=TH_0=q+R,\qquad R=r+kK^3g,\qquad 0\le B\le2.
\tag{Q62.2.2}\label{eq:Q62:Q62:2:2}
\]


Expanding (Q62.2.1) gives

\[
\begin{aligned}
TP&\le a(q+r)+tK^3g-tK(gs)-t\delta_2+\tfrac14T(qs^2)\\
&\le aB+t\bigl(|K(gs)|-\delta_2\bigr)+\tfrac12E(qs^2).
\end{aligned}
\]

Here $T(qs)=\allowbreak \delta_2-K^3g+K(gs)$. The row bound gives
$|K(gs)|\le\allowbreak \sqrt{d\delta_3}\le\allowbreak \sqrt{\theta d\delta_2}$.
For $d>\allowbreak 0$, put $v=\allowbreak \delta_2/d$. Completing the square in
$\sqrt{\theta v}-v\le\allowbreak \theta/4$ retains the negative term $-t\delta_2$.
Also $E(qs^2)/2\le\allowbreak  \delta_3\le\allowbreak \theta^2d$. Hence

\[
0\le TP\le t^2(q+R)+\gamma d,
\qquad\gamma=t\theta/4+\theta^2.
\tag{Q62.2.3}\label{eq:Q62:Q62:2:3}
\]


Taking the arithmetic mean of the nonnegative quantities $tp$ and $t(q+R)+\gamma d/t$ yields

\[
0\le pTP\le(c+N/2)^2,
\quad N=tR=K(tg+s),
\quad c=t+\gamma d/(2t).
\tag{Q62.2.4}\label{eq:Q62:Q62:2:4}
\]

The base of this square is nonnegative, being their arithmetic mean.
As $t\ge\allowbreak 5/6$,

\[
c\le1-\frac{1793}{5120}d\le1-\frac7{20}d=:c_0.
\tag{Q62.2.5}\label{eq:Q62:Q62:2:5}
\]

The last comparison has slack $d/5120$. Enlarging the base preserves its nonnegativity.

Retain the same parameter $v=\allowbreak \delta_2/d$ throughout. Then

\[
d\le v\le\theta,\qquad m=t-v>0,\qquad EN=-m d.
\tag{Q62.2.6}\label{eq:Q62:Q62:2:6}
\]

The spectral measure $\sigma$ of $g$ has mass $d$, support $[0,\allowbreak  \theta]$,
and first moment $\delta_2$. The second derivative of $x(t-x)^2$ is $-4t+6x<\allowbreak 0$,
since $-4(5/6)+6(13/32)<\allowbreak 0$. Jensen gives

\[
EN^2=\int x(t-x)^2\,d\sigma(x)
\le v(t-v)^2d.
\tag{Q62.2.7}\label{eq:Q62:Q62:2:7}
\]

Here Jensen is applied to the probability measure $\sigma/d$, with mean $v$.
The row bound and $\delta_3\le\allowbreak \theta \delta_2$ also give

\[
\begin{gathered}
|N|^2\le\|tg+s\|_2^2\\
=t^2d-2t\delta_2+\delta_3\\
\le d\{t^2-(2t-\theta)v\}=:U^2.
\end{gathered}\tag{Q62.2.8}\label{eq:Q62:Q62:2:8}
\]

Since $2t-\theta>\allowbreak 0$ and $v\ge\allowbreak  d$,

\[
U^2\le f(d):=d-\frac{83}{32}d^2+\frac54d^3<\frac19.
\tag{Q62.2.9}\label{eq:Q62:Q62:2:9}
\]

To verify the last comparison, first take $d\le\allowbreak 1/4$. Then
$f(d)\le\allowbreak  d-(73/32)d^2\le\allowbreak 8/73<\allowbreak 1/9$.
For $1/4\le\allowbreak  d\le\allowbreak 1/3$,
$f^{\prime\prime}(d)=\allowbreak -83/16+(15/2)d\le\allowbreak -43/16<\allowbreak 0$ and $f^\prime(1/4)=\allowbreak -1/16$.
Thus $f$ decreases on $[1/4,\allowbreak  1/3]$ and has maximum $f(1/4)=\allowbreak 55/512<\allowbreak 1/9$.
Take $U$ to be the nonnegative square root; at zero use continuity.

\paragraph{Q62.3 A scalar exponential bound}

For $U\ge\allowbreak 0$, define

\[
C(U)=
\begin{cases}(e^{8U}-1-8U)/U^2,&U>0,\\32,&U=0.
\end{cases}
\]

The integral form of the Taylor remainder gives
$C(U)=\allowbreak 64\int_0^1(1-z)e^{8zU}\,\allowbreak  dz$.
This representation holds for every real argument and is increasing in that argument. Consequently, for $n\le\allowbreak  U$,

\[
e^{8n}\le1+8n+C(U)n^2.
\tag{Q62.3.1}\label{eq:Q62:Q62:3:1}
\]

For $0\le\allowbreak  U\le\allowbreak 1/3$, we claim

\[
\boxed{C(U)\le195/2,\qquad C(U)\le48+448U^2.}
\tag{Q62.3.2}\label{eq:Q62:Q62:3:2}
\]

The terms of orders zero through four in $e^{4/3}$ sum to $911/243$.
The fifth term is $128/3645$, and every subsequent ratio is at most $2/9$. Therefore

\[
\begin{gathered}
e^{4/3}\le\frac{911}{243}+\frac{128}{3645}\frac9{7}\\
=\frac{32269}{8505}<\frac{19}{5}.
\end{gathered}
\]

The final numerator difference is 50. It follows that
$e^{8/3}<\allowbreak (19/5)^2=\allowbreak 361/25<\allowbreak 29/2$,
so $C(1/3)<\allowbreak 195/2$; monotonicity proves the first bound.

Separate the first two terms of $C(U)$. The remaining coefficients are nonnegative and their degrees are at least two. For $U\le\allowbreak 1/3$,

\[
C(U)\le32+\frac{256}{3}U+\frac{667}{2}U^2.
\]

The quadratic coefficient is
$9\{195/2-32-256/9\}=\allowbreak 667/2$.
The nonnegative square $(32U/3-4)^2$ gives

\[
\frac{256}{3}U\le16+\frac{1024}{9}U^2.
\]

Hence $C(U)\le\allowbreak 48+(8051/18)U^2\le\allowbreak 48+448U^2$, proving (Q62.3.2).


\paragraph{Q62.4 Reduction to two parameters of the same kernel}

By (Q62.2.4)--(Q62.2.5) and $1+x\le\allowbreak  e^x$,

\[
E(pTP)^8\le E(c_0+N/2)^{16}
\le e^{-28d/5}Ee^{8N}.
\]

Now apply (Q62.3.1), (Q62.2.6)--(Q62.2.8), and $1+x\le\allowbreak  e^x$:

\[
\boxed{\begin{gathered}
E(pTP)^8\le\exp\{d\,\mathcal E(d,v)\},\quad\\
\mathcal E(d,v)=-\frac{28}{5}-8(t-v)+C(U)v(t-v)^2.
\end{gathered}}\tag{Q62.4.1}\label{eq:Q62:Q62:4:1}
\]

The integrated polynomial upper bound is positive because it dominates $Ee^{8N}>\allowbreak 0$.
It can therefore be exponentiated.
We prove, for $0<\allowbreak d\le\allowbreak 1/3$ and $d\le\allowbreak  v\le\allowbreak 13/32$, that

\[
\mathcal E(d,v)\le-1/200.
\tag{Q62.4.2}\label{eq:Q62:Q62:4:2}
\]

\paragraph{Q62.5 The range $v\le\allowbreak 1/4$}

Here $d\le\allowbreak 1/4$. Since $t\ge\allowbreak 7/8$, the function $v(t-v)^2$ increases on $[0,\allowbreak  1/4]$, and

\[
\begin{gathered}
v(t-v)^2\le\frac14(t-1/4)^2\\
=\frac9{64}-\frac3{16}d+\frac1{16}d^2,\\
\quad -\frac{28}{5}-8(t-v)\le-\frac{58}{5}+4d.
\end{gathered}\tag{Q62.5.1}\label{eq:Q62:Q62:5:1}
\]

\paragraph{Q62.5.1 The range $0<\allowbreak d\le\allowbreak 3/16$}

Using $d^3\le\allowbreak (3/16)d^2$ in (Q62.2.9) gives

\[
U^2\le d-\frac{151}{64}d^2,
\qquad C(U)\le48+448d-1057d^2.
\]

Multiplying by the bounds in (Q62.5.1) and expanding gives

\[
\begin{gathered}
\mathcal E\le-\frac{97}{20}+58d-\frac{14697}{64}d^2\\
+\frac{3619}{16}d^3-\frac{1057}{16}d^4.
\end{gathered}
\]

Drop the final nonpositive term and use $d^3\le\allowbreak (3/16)d^2$ to obtain

\[
\begin{gathered}
\mathcal E\le-\frac{97}{20}+58d-\frac{47931}{256}d^2\\
\le-\frac{97}{20}+\frac{215296}{47931}<-\frac14.
\end{gathered}\tag{Q62.5.2}\label{eq:Q62:Q62:5:2}
\]

The last comparison follows from $5\cdot215296=\allowbreak 1076480<\allowbreak 23\cdot47931=\allowbreak 1102413$.
The quadratic upper bound is obtained by completing the square and holds for every real $d$.

\paragraph{Q62.5.2 The range $3/16\le\allowbreak  d\le\allowbreak 1/4$}

The variance upper bound in (Q62.5.1) decreases with $d$, so it is at most $441/4096$.
Also $C(U)\le\allowbreak 195/2$ and $-58/5+4d\le\allowbreak -53/5$. Thus

\[
\begin{gathered}
\mathcal E\le\frac{195}{2}\frac{441}{4096}-\frac{53}{5}\\
=-\frac{4201}{40960}<-\frac1{10}.
\end{gathered}\tag{Q62.5.3}\label{eq:Q62:Q62:5:3}
\]



\paragraph{Q62.6 The range $1/4\le\allowbreak  v\le\allowbreak 5/16$}

Now $d\le\allowbreak 5/16$. On $0\le\allowbreak  v\le\allowbreak  t$, the maximum of $v(t-v)^2$ is $4t^3/27$.
Moreover,

\[
-\frac{28}{5}-8(t-v)\le-\frac{111}{10}+4d.
\]

Using $v\ge\allowbreak 1/4$ in (Q62.2.8) gives

\[
\begin{gathered}
U^2\le d\left(\frac{77}{128}-\frac34d+\frac14d^2\right)\\
\le\frac{77}{128}d-\frac{43}{64}d^2,\\
\quad C(U)\le48+\frac{539}{2}d-301d^2.
\end{gathered}\tag{Q62.6.1}\label{eq:Q62:Q62:6:1}
\]

The second step uses $d\le\allowbreak 5/16$. For $0\le\allowbreak  d\le\allowbreak 1/3$,

\[
t^3=(1-d/2)^3\le1-\frac54d,
\tag{Q62.6.2}\label{eq:Q62:Q62:6:2}
\]

because $(3/4)d^2\le\allowbreak  d/4$ and the remaining cubic term is nonpositive.
The bound in (Q62.6.1) dominates the positive quantity $C(U)$ and can be multiplied by (Q62.6.2). This yields

\[
\begin{aligned}
\mathcal E
&\le-\frac{359}{90}+\frac{946}{27}d-\frac{189}{2}d^2
+\frac{1505}{27}d^3\\
&\le-\frac{359}{90}+\frac{946}{27}d-\frac{33299}{432}d^2\\
&\le-\frac{359}{90}+\frac{3579664}{899073}<-\frac1{200}.
\end{aligned}
\tag{Q62.6.3}\label{eq:Q62:Q62:6:3}
\]

The second line uses $d^3\le\allowbreak (5/16)d^2$; the third completes the square.
The final rational comparison follows from

\[
\begin{gathered}
3579664\cdot1800=6443395200\\
<899073\cdot7171=6447252483,\\
\quad \frac{359}{90}-\frac1{200}=\frac{7171}{1800}.
\end{gathered}
\]



\paragraph{Q62.7 The range $5/16\le\allowbreak  v\le\allowbreak \theta$}

By (Q62.2.8) and $d\le\allowbreak 1/3$,

\[
\begin{aligned}
U^2&\le d\left(\frac{257}{512}-\frac{11}{16}d+\frac14d^2\right)
\le\frac{257}{512}d-\frac{29}{48}d^2,\\
C(U)&\le48+\frac{1799}{8}d-\frac{812}{3}d^2.
\end{aligned}
\tag{Q62.7.1}\label{eq:Q62:Q62:7:1}
\]

For the mean term, use $v\le\allowbreak \theta$ to obtain
$-28/5-8(t-v)\le\allowbreak -207/20+4d$.
For the variance use $4t^3/27$ and (Q62.6.2). Expanding and absorbing the cubic term gives

\[
\begin{aligned}
\mathcal E
&\le-\frac{583}{180}+\frac{1535}{54}d
-\frac{52969}{648}d^2+\frac{4060}{81}d^3\\
&\le-\frac{583}{180}+\frac{1535}{54}d-\frac{126427}{1944}d^2\\
&\le-\frac{583}{180}+\frac{4712450}{1517124}<-\frac1{10}.
\end{aligned}
\tag{Q62.7.2}\label{eq:Q62:Q62:7:2}
\]

The second line uses $d^3\le\allowbreak  d^2/3$; the third completes the square.
The last comparison follows from

\[
\begin{gathered}
4712450\cdot36=169648200\\
<1517124\cdot113=171435012,\\
\quad \frac{583}{180}-\frac1{10}=\frac{113}{36}
\end{gathered}
\]

The bounds (Q62.5.2), (Q62.5.3), (Q62.6.3), and (Q62.7.2) cover the allowed parameter domain.
Together they prove (Q62.4.2), and hence $EH^8\le\allowbreak  e^{-d/200}$.
Finally,

\[
e^{-d/200}\le\frac1{1+d/200}
=1-\frac{d}{200+d}\le1-\frac{d}{601}
\]

since $200+d\le\allowbreak 601$.


\paragraph{Q62.8 Boundaries, decorations and reversal}

If $d=\allowbreak 0$, then $g=\allowbreak 0$ and $A\one=\allowbreak T\one=\allowbreak \one$. Leaf integration gives every path message equal to one,
so the conclusion holds with equality. Divisions by $v$ or $U$ above are taken for positive parameters or by the stated continuous extension.
Omit zero-mass atoms. On the remaining support, $p,\allowbreak  q\ge\allowbreak \mu(x)>\allowbreak 0$.


As $E_0\le\allowbreak 1$ and $T$ is positive, the decorated message is bounded by the undecorated one.
Equal-arm factors at the receiving root can be restored in the same way.
Replacing $K$ by $-K$ reverses all edges and gives the bound for $qAF^-$ with the same $d$.
